\documentclass[10pt,a4paper]{amsart}
\usepackage[margin=19mm]{geometry}
\usepackage{amsmath,amssymb,mathtools}
\usepackage[scr=rsfs,cal=euler]{mathalfa}
\usepackage{xfrac}
\usepackage[unicode,hidelinks,bookmarks=false]{hyperref}
\newcommand{\ie}{i.e.\ }
\newcommand{\cf}{cf.\ }
\newcommand{\resp}{respectively\ }
\newcommand{\Subsection}{\subsection}
\providecommand{\nobreakdash}{\nobreak\hbox{-}}
\setlength{\emergencystretch}{2em}
\multlinegap0pt
\usepackage{bm}
\usepackage{bbm}
\usepackage{dsfont}
\usepackage{xspace}
\usepackage{cancel}
\usepackage{mathtools}
\usepackage{amsthm}
\makeatletter
\@ifundefined{theorem}{\newtheorem{theorem}{Theorem}}{}
\makeatother
\title[Spencer-Riemann-Roch Theory: Mirror Symmetry of Hodge Decomp]{Spencer-Riemann-Roch Theory: Mirror Symmetry of Hodge Decompositions and Characteristic Classes in Constrained Geometry}
\author{Dongzhe Zheng}
\date{Source: arXiv:2506.05915v2 (CC BY 4.0)}
\begin{document}
\maketitle

\noindent\emph{Source and licence.} Spencer-Riemann-Roch Theory: Mirror Symmetry of Hodge Decompositions and Characteristic Classes in Constrained Geometry. Version arXiv:2506.05915v2 (CC BY 4.0), distributed under the Creative Commons Attribution 4.0 International licence, as recorded in the arXiv licence metadata for this exact version and shown on the abstract page https://arxiv.org/abs/2506.05915v2. Full rights notice: licenses/arxiv-2506.05915-CC-BY-4.0.txt.\par\smallskip

\noindent\textit{Editorial scope.} Counted unit: the Theorem whose source label is \texttt{eq:alg-hodge-detailed} in the article (source0.tex lines 253--259, source label \texttt{eq:alg-hodge-detailed}), delivered together with the article's own complete original proof of it (source0.tex lines 261--298, one \texttt{proof} environment of 38 lines). No mathematics is written, repaired or re-derived: statement and proof are the article's own text line for line. Required context retained uncounted: none. Every definition, display and label of those lines is retained unchanged. Omitted from this package and not redistributed: 1--252, 260--260, 299--939. Nothing inside a retained excerpt is omitted or altered: every retained body is delivered exactly as the article's lines are written, so no omission receipt of any kind accompanies this package. Citations: the retained excerpts carry no citation command at all, so no bibliography entry of the article is transcribed and no reference is redistributed. Reference adaptation: none was needed and none was made; every reference inside the delivered unit resolves inside it, so no unresolved reference or citation is produced. Printed slips kept literal: no printed word, symbol, hypothesis or number of the counted statement or of its proof was altered. Layout and class: the article's own class is \texttt{extarticle} with packages including inputenc, geometry, graphicx, booktabs, mathrsfs, bm, xcolor, algorithm, algpseudocode, mathtools, enumitem, tikz-cd, newtxtext, newtxmath; the wrapper is the lane's pinned \texttt{amsart} editorial wrapper at 10pt on A4, which restates the theorem-like environments the retained text needs and adds the article's own macro definitions that the retained text needs (none), with extra wrapper packages (bm, bbm, dsfont, xspace, cancel, mathtools). The delivered numbering is the unit's own. Assets: no retained span contains a figure environment, a graphics-inclusion command, a PDF-page inclusion, a table or a TikZ picture; no image file is copied into this package, the package declares no asset dependency and no image file is redistributed with it. Licence: version arXiv:2506.05915v2 of this article is distributed under the Creative Commons Attribution 4.0 International licence, as recorded in the arXiv licence metadata for this exact version and shown on the abstract page https://arxiv.org/abs/2506.05915v2. A full rights notice is delivered at licenses/arxiv-2506.05915-CC-BY-4.0.txt. Characters: every retained excerpt is pure ASCII, so the delivered engine, XeLaTeX with its default Latin Modern text font, reports no missing character.
\begin{theorem}[Algebraic Geometric Version of Spencer-Hodge Decomposition]
Let $(D,\lambda)$ be a compatible pair satisfying strong transversality conditions, with Lie group $G$ satisfying strict conditions. Then there exists an orthogonal decomposition:
\begin{equation}\label{eq:alg-hodge-detailed}
\mathcal{S}^k_{D,\lambda} = \mathcal{H}^k_{D,\lambda} \oplus \mathcal{D}^{k-1}(\mathcal{S}^{k-1}_{D,\lambda}) \oplus \mathcal{D}^{k+1*}(\mathcal{S}^{k+1}_{D,\lambda})
\end{equation}
where $\mathcal{H}^k_{D,\lambda} := \ker(\Delta^k_{D,\lambda}) \cap \mathcal{S}^k_{D,\lambda}$ is the algebraic geometric harmonic space.
\end{theorem}

\begin{proof}
The proof is based on standard theory of elliptic operators, with particularly strong forms under strict Lie group conditions:

\textbf{Step 1: Ellipticity verification}. Strong transversality conditions ensure $\Delta^k_{D,\lambda}$ is an elliptic operator. Specifically, the principal symbol of Spencer differential $\mathcal{D}^k_{D,\lambda}$ is determined by the principal symbol of exterior differential $d$:
$$\sigma(\mathcal{D}^k_{D,\lambda})(\xi) = \sigma(d)(\xi) \otimes \text{id}_{\mathrm{Sym}^k(\mathcal{G})}$$
where $\xi \in T^*M$ is a non-zero cotangent vector. The principal symbol $\sigma(d)(\xi): \Omega^k_M \to \Omega^{k+1}_M$ of exterior differential $d$ is given by $\xi \wedge (\cdot)$, which is always injective for $\xi \neq 0$. Strong transversality conditions ensure transmission of ellipticity by ensuring non-degeneracy of compatible pair structures, well-definedness and non-degeneracy of Spencer extension operators $\delta^{\lambda}_{\mathcal{G}}$, maintaining invertibility properties of principal symbols using trivial center, and ensuring standard results of elliptic operator theory apply through compactness. Under strict conditions, ellipticity is particularly stable.

\textbf{Step 2: Finite dimensionality}. Harmonic spaces of elliptic operators on compact manifolds are finite dimensional:
$$\dim \mathcal{H}^k_{D,\lambda} < \infty$$
This is guaranteed by standard theory of elliptic operators. Strict Lie group conditions ensure stability and computability of this dimension.

\textbf{Step 3: Existence of Green operators}. Ellipticity guarantees existence of Green operators $G^k: \mathcal{S}^k_{D,\lambda} \to \mathcal{S}^k_{D,\lambda}$ satisfying:
$$\Delta^k_{D,\lambda} G^k = \text{id} - P^k$$
where $P^k: \mathcal{S}^k_{D,\lambda} \to \mathcal{H}^k_{D,\lambda}$ is harmonic projection. Construction of Green operators is based on fundamental solution theory of elliptic operators, with existence and uniqueness guaranteed by ellipticity and compactness. Under strict conditions, Green operators have particularly good regularity and algebraic properties.

\textbf{Step 4: Construction of orthogonal decomposition}. Use Green operators to construct orthogonal decomposition. First, harmonic projection $P^k$ decomposes $\mathcal{S}^k_{D,\lambda}$ into harmonic and non-harmonic parts:
$$\mathcal{S}^k_{D,\lambda} = \mathcal{H}^k_{D,\lambda} \oplus (\mathcal{H}^k_{D,\lambda})^{\perp}$$
Then, for the non-harmonic part, use properties of Green operators: for any $\omega \in (\mathcal{H}^k_{D,\lambda})^{\perp}$, we have
$$\omega = \mathcal{D}^{k-1}((\mathcal{D}^{k-1})^* G^k \omega) + (\mathcal{D}^{k+1})^*(\mathcal{D}^{k+1} G^k \omega)$$
This establishes the decomposition:
$$(\mathcal{H}^k_{D,\lambda})^{\perp} = \mathrm{im}(\mathcal{D}^{k-1}) \oplus \mathrm{im}(\mathcal{D}^{k+1*})$$

\textbf{Step 5: Verification of orthogonality}. Verify orthogonality between subspaces by direct computation:

For $\langle \mathcal{H}^k_{D,\lambda}, \mathrm{im}(\mathcal{D}^{k-1}) \rangle = 0$: Let $h \in \mathcal{H}^k_{D,\lambda}$ and $\mathcal{D}^{k-1}\alpha \in \mathrm{im}(\mathcal{D}^{k-1})$, then
$$\langle h, \mathcal{D}^{k-1}\alpha \rangle = \langle \mathcal{D}^{k-1*}h, \alpha \rangle = 0$$
because $h$ is harmonic, so $\mathcal{D}^{k-1*}h = 0$.

For $\langle \mathcal{H}^k_{D,\lambda}, \mathrm{im}(\mathcal{D}^{k+1*}) \rangle = 0$: Let $h \in \mathcal{H}^k_{D,\lambda}$ and $\mathcal{D}^{k+1*}\beta \in \mathrm{im}(\mathcal{D}^{k+1*})$, then
$$\langle h, \mathcal{D}^{k+1*}\beta \rangle = \langle \mathcal{D}^{k+1}h, \beta \rangle = 0$$
because $h$ is harmonic, so $\mathcal{D}^{k+1}h = 0$.

For $\langle \mathrm{im}(\mathcal{D}^{k-1}), \mathrm{im}(\mathcal{D}^{k+1*}) \rangle = 0$: Let $\mathcal{D}^{k-1}\alpha \in \mathrm{im}(\mathcal{D}^{k-1})$ and $\mathcal{D}^{k+1*}\beta \in \mathrm{im}(\mathcal{D}^{k+1*})$, then
$$\langle \mathcal{D}^{k-1}\alpha, \mathcal{D}^{k+1*}\beta \rangle = \langle \mathcal{D}^{k+1}\mathcal{D}^{k-1}\alpha, \beta \rangle = 0$$
because $\mathcal{D}^{k+1}\mathcal{D}^{k-1} = 0$ (property of Spencer complexes).

\textbf{Step 6: Special advantages of strict conditions}. Semi-simplicity ensures non-degeneracy of Hodge decomposition, as non-degeneracy of Killing forms transmits to all relevant bilinear forms; trivial center avoids unnecessary reduction steps, simplifying analysis and ensuring uniqueness of results; compactness ensures convergence of all infinite series, particularly integral representations of Green operators; algebraicity makes decomposition well-defined in the algebraic geometric category, consistent with analytic category results through the GAGA principle.
\end{proof}

\end{document}
